\section{Conclusion}
\label{sec:conclusion}

In this work, we provide a full calibration of the four HSC tomographic bins, including constraints on the two higher redshift tomographic bins. 
With the arrival of Stage-IV weak lensing surveys such as LSST, Euclid and the Roman Space Telescope, and further spectroscopic data programs available such as the next data releases of DESI and the data acquisition of the Prime Focus Spectrograph (PFS) \cite{2016PFSspectro} program, systematic mitigation will become significantly more important. 
Therefore, this work explored reducing some of the common systematics in clustering redshifts, such as galaxy bias and magnification effects. 
The fiducial scale cut choice of this analysis is $0.3-3\,\hMpc$. Consistency between the two presented scale cuts ($0.3-3\,\hMpc$ and $1-5\,\hMpc$) alleviates concerns that bias modeling complexities have substantial impacts on the showcased results in section \ref{sec:results}.

% Should we mention the companion work here ? YES
%This work comes with a companion paper \cite{ChoppinDeJanvry2025b} where cosmological parameter inference on cosmic shear in both real and Fourier space is ran with the calibrated data products of this analysis and results are compared in the $S_8-\Omega_m$ plane with the HSC Year 3 cosmic shear measurements \cite{CosmoShearLi2023HSCY3,Dalal2023CosmoShearHSC} and other results from weak lensing surveys (KiDS-Legacy \cite{wright2025kidslegacyCosmoShear}, DES \cite{Secco2022CosmoShearDES,Amon2022CosmoShearDES}) as well as Planck \cite{Planck18}. 

Our calibration of the tomographic distributions suggests smaller shifts than expected by the cosmic shear analyses from HSC Year 3 and the shear calibration ratios analysis in Bin 3 and Bin 4. 
This analysis also hints at biases present in Bin 1 and a mild bias present in Bin 2, both towards lower redshifts. 
In companion paper \cite{ChoppinDeJanvry2025b}, we present the cosmological parameter inference on cosmic shear in Fourier space with the calibrated redshift distributions from this analysis.

%Our calibration of the tomographic distributions suggests smaller shifts than expected by the cosmic shear analyses from HSC Year 3 and the shear calibration ratios analysis: as such, we expect the $S_8$ constraints to shift towards higher values, possibly obtaining better consistency of HSC with Planck. 
%in the plane $S_8-\Omega_m$, 
%For reference, the $S_8$ values obtained by the real space two-point correlation functions are $S_{8}=0.769^{+0.031}_{-0.034} \; (0.782)$ \cite{CosmoShearLi2023HSCY3} and $S_8=0.776^{+0.032}_{-0.033} \; (0.792)$ for the angular power spectrum \cite{Dalal2023CosmoShearHSC}, with MAP values in parentheses, while Planck reports $S_{8}=0.830 ± 0.013$ in Table 1 of \cite{Planck18}. The calibrated distributions made in this work offer new insights on constraining the growth of structure with HSC and path the way towards more clustering redshifts re-analyses of the major existing and upcoming weak lensing datasets. We especially improve constraints in the higher redshift ranges, with DESI's ELG and QSO target classes. Further improvements can definitely also be expected with the next data releases of the DESI experiment, which we leave to future works.
The code used for this work and all the data to replicate the figures of this paper is made publicly available at \href{https://github.com/JeanCHDJdev/desi-y3-hsc/}{\textcolor{blue}{github.com/jeanchdjdev/desi-y3-hsc}}.